% =====================================================================
%  Rotinas de quem? Inspeção semiótica das Rotinas da Alexa
%  Template: SBC acessível de uma coluna, usado no IHC 2026 (Quixadá/CE)
%  Compilar: latexmk -pdf main.tex   (ou pdflatex + bibtex + 2x pdflatex)
% =====================================================================
\documentclass[12pt]{article}

\usepackage{sbc-template}   % A4, Times 12 pt, margens 3,5/2,5/3/3 cm

\usepackage[brazilian]{babel}
\usepackage[T1]{fontenc}
\usepackage[utf8]{inputenc}
\usepackage{graphicx,url}
\usepackage{multirow}
\usepackage{array}
\usepackage{textcomp}
\usepackage{soul}
\usepackage[table,xcdraw]{xcolor}
\usepackage[shortlabels]{enumitem}
\usepackage{comment}
\usepackage{fancyhdr}
\usepackage{IHC}            % cabeçalho com evento e trilha
\usepackage[super]{nth}
\usepackage{pdfcomment}     % descrição de imagens para leitores de tela (\pdftooltip)

% Configurações de hifenização do template oficial
\tolerance=1
\emergencystretch=\maxdimen
\hyphenpenalty=10000
\hbadness=10000
\sloppy

\graphicspath{{figures/}}

% O pdfcomment carrega o hyperref, que por padrão desenha caixas coloridas
% em volta de citações e links. hidelinks mantém os links clicáveis sem caixas.
\hypersetup{hidelinks}

% ---------------------------------------------------------------------
%  Citação no texto corrido, no padrão SBC: \citet{chave} gera
%  "Geeng e Roesner [2019]". O \cite{chave} continua gerando
%  "[Geeng e Roesner 2019]". O sbc.bst não oferece esse comando.
% ---------------------------------------------------------------------
\ExplSyntaxOn
\tl_new:N \l__artigo_cite_tl
\NewDocumentCommand{\citet}{m}
  {
    \nocite{#1}
    \cs_if_exist:cTF { b@#1 }
      {
        \tl_set_eq:Nc \l__artigo_cite_tl { b@#1 }
        % Com o hyperref, o rótulo fica dentro de um link: o ano é
        % seguido de um fecha-chaves, capturado em \2 e devolvido.
        \regex_replace_once:nnN { \s (\d{4}[a-z]?) (\cE.)? \Z } { \ [\1]\2 } \l__artigo_cite_tl
        \tl_use:N \l__artigo_cite_tl
      }
      { \textbf{?} }
  }
\ExplSyntaxOff

% ---------------------------------------------------------------------
%  Marcação de trechos pendentes. Troque para \pendentesfalse para
%  gerar a versão sem as marcações em vermelho.
% ---------------------------------------------------------------------
\newif\ifpendentes
\pendentestrue
\newcommand{\pendente}[1]{\ifpendentes{\color{red}[PENDENTE: #1]}\fi}

% Dados do evento exibidos no cabeçalho (template IHC 2026)
\trilha{Pesquisa} %Preencha com: Pesquisa | Ideias Inovadoras e Resultados Emergentes | Relatos de Experiência | Prêmio Junia Coutinho Anacleto | Pôsteres, Demonstrações e Artigos Internacionais | WTD ! Minicursos e Workshops | Competição de Avaliação | IHC na Prática | WEIHC | WAIHCWS
\local{Quixadá/CE}
\ano{2026}
\edicao{XXV}


% ---------------------------------------------------------------------
%  Título e autoria (o título deve estar no idioma do texto)
% ---------------------------------------------------------------------
\title{Rotinas de quem? Inspeção semiótica das Rotinas da Alexa\\
sob a perspectiva da casa como usuário coletivo}

\author{Nome Sobrenome\inst{1}, Nome Sobrenome\inst{1}}

\address{Departamento -- Universidade\\
  Cidade -- UF -- Brasil
  \email{autor@exemplo.br}
}

\begin{document}

\maketitle
\thispagestyle{plain}

% Resumo estruturado em inglês e em português (obrigatório no IHC)
% =====================================================================
%  Resumo (LNCS: 150 a 250 palavras) e palavras-chave
% =====================================================================
\begin{abstract}
As Rotinas da Amazon Alexa permitem associar um evento a uma sequência de
ações executadas pelo assistente e pelos dispositivos conectados. Elas são
criadas por uma pessoa, no aplicativo ligado à sua conta, e executadas em
espaços compartilhados da casa, onde afetam todos que estão presentes. Este
artigo avalia a metacomunicação do aplicativo Amazon Alexa na criação e no
gerenciamento de rotinas, considerando a casa como usuário coletivo. Para
isso, aplicamos o Método de Inspeção Semiótica com uma lente para múltiplos
interlocutores, que distingue o que a interface comunica a quem cria a
rotina e o que comunica sobre quem é afetado por ela. Os achados são
triangulados com avaliações públicas do aplicativo nas lojas brasileiras.
\pendente{resumir os principais resultados e implicações de design após a
inspeção}
\keywords{Engenharia Semiótica \and Método de Inspeção Semiótica \and
Assistentes de voz \and Casas inteligentes \and Rotinas \and IHC}
\end{abstract}


% ---------------------------------------------------------------------
%  Corpo do artigo (uma seção por arquivo)
% ---------------------------------------------------------------------
% =====================================================================
%  1 Introdução
% =====================================================================
\section{Introdução}
\label{sec:introducao}

Assistentes de voz instalados em alto-falantes inteligentes, como a Amazon
Alexa nos dispositivos Echo, passaram a fazer parte do cotidiano doméstico.
Estudos baseados em registros de uso mostram que ouvir música, fazer buscas
e controlar outros dispositivos da casa estão entre as atividades mais
frequentes com esses assistentes~\cite{ammari2019music}. Diferentemente de
um smartphone, o alto-falante costuma ficar em um espaço comum da casa, e
seu uso se entrelaça com as conversas e as atividades simultâneas dos
moradores, como as refeições em família~\cite{porcheron2018voice}.

Entre os recursos oferecidos pela Alexa estão as Rotinas, que permitem
associar um evento disparador (um horário, um comando de voz, um alarme) a
uma sequência de ações executadas pelo assistente e pelos dispositivos
conectados. A criação de rotinas segue o modelo de regras gatilho-ação
típico do desenvolvimento pelo usuário final e é feita, principalmente, por
meio do aplicativo complementar instalado no
smartphone~\cite{barricelli2024routines,barricelli2022multimodal}. As
pesquisas sobre o tema concentram-se em facilitar essa autoria, comparando
a criação pelo aplicativo, por conversa e por interação multimodal em
alto-falantes com tela~\cite{barricelli2024routines,barricelli2022multimodal}.

Há, porém, uma diferença entre quem cria uma rotina e quem a vivencia. A
rotina é configurada por uma pessoa, no aplicativo associado à sua conta, e
executada em um ambiente compartilhado, onde afeta todos que estão
presentes. Estudos sobre casas inteligentes com múltiplos usuários mostram
que a instalação e a configuração tendem a se concentrar em um morador, o
que gera assimetrias de controle e tensões com os demais
residentes~\cite{geeng2019control,zeng2019multiuser}. Pessoas que trabalham
ou convivem no espaço doméstico, como babás, também são afetadas por esses
sistemas sem participar das decisões sobre eles~\cite{bernd2020bystanders}.
Sob a perspectiva da integridade contextual~\cite{nissenbaum2004privacy},
quando um dado ligado a uma pessoa, como seu horário ou sua voz, dispara
uma ação em um espaço coletivo, a informação passa a circular em um
contexto diferente daquele em que foi fornecida.

Os trabalhos citados investigam essas questões empiricamente, com
moradores, ou avaliam a usabilidade da autoria de rotinas considerando um
único autor. Recebe menos atenção a forma como a própria interface comunica
o que é uma rotina, quem a controla e a quem ela afeta. A Engenharia
Semiótica oferece uma lente para essa análise ao tratar a interface como
uma mensagem do designer para o usuário, isto é, como
metacomunicação~\cite{desouza2005semiotic}. O Método de Inspeção Semiótica
(MIS), derivado dessa teoria, permite reconstruir essa mensagem a partir
dos signos da interface, sem a participação de
usuários~\cite{desouza2006sim,desouza2009methods}, e já foi adaptado a
diferentes contextos de aplicação~\cite{fortunato2025ihc,fortunato2025adaptations}. Até onde
pudemos verificar, o método ainda não foi aplicado à criação e ao
gerenciamento de rotinas em assistentes de voz. Além disso, a MIS pressupõe
um destinatário da metamensagem, enquanto no caso das rotinas há ao menos
dois: quem cria a rotina e quem é afetado por ela.

Essa lacuna se relaciona com os Grandes Desafios de Pesquisa em IHC no
Brasil para 2025-2035 (GranDIHC-BR)~\cite{pereira2024grandihc}. O trabalho
se insere principalmente no GC4, que trata dos aspectos socioculturais em
IHC~\cite{neris2024gc4}, ao examinar uma tecnologia concebida em torno de
uma conta individual e usada em arranjos domésticos coletivos. Contribui
também para o GC1, voltado a novas abordagens teóricas e
metodológicas~\cite{dasilvajunior2024gc1}, ao propor uma adaptação da MIS
que considera múltiplos interlocutores. Por fim, dialoga com o GC7, que
aborda a interação com tecnologias emergentes como um ecossistema que
integra pessoas, tecnologias e contextos~\cite{zaina2024gc7}.

Diante disso, este artigo tem como objetivo avaliar a metacomunicação do
aplicativo Amazon Alexa na criação e no gerenciamento de rotinas,
considerando a casa como usuário coletivo, e derivar implicações de design
para lares compartilhados. A pesquisa é orientada pela seguinte questão
principal: como o aplicativo Amazon Alexa comunica a criação e o
funcionamento de rotinas, e a quem essa comunicação se dirige em uma casa
compartilhada? Dela derivam duas questões específicas:

\begin{itemize}
  \item \textbf{QP1:} Que metamensagem o designer transmite sobre autoria,
        uso de dados e pessoas afetadas pelas rotinas?
  \item \textbf{QP2:} Que potenciais rupturas de comunicação a inspeção
        identifica para a persona que administra a conta e para a persona
        que convive com as rotinas sem administrá-las?
\end{itemize}

Para responder a essas questões, aplicamos a MIS ao aplicativo Amazon
Alexa em português. A inspeção foi guiada por duas personas construídas
a partir da literatura sobre casas compartilhadas, a pessoa que
administra a conta e cria as rotinas e um morador que convive com elas
sem administrá-las, e por cenários de uso para cada uma. Em cada etapa,
os signos foram analisados quanto ao que comunicam a cada persona.
Nenhum participante foi recrutado: as personas são fictícias e a
inspeção usou contas dos próprios pesquisadores.

As contribuições deste trabalho são:

\begin{itemize}
  \item a reconstrução da metamensagem do aplicativo de Rotinas da Alexa,
        evidenciando a quem o design se dirige;
  \item um catálogo de potenciais rupturas de comunicação para quem
        administra a conta e para quem convive com as rotinas;
  \item uma adaptação da MIS que inclui uma lente para múltiplos
        interlocutores, aplicável a outros sistemas de uso coletivo;
  \item duas personas fundamentadas na literatura, com a origem de cada
        característica documentada, que podem ser reutilizadas em outros
        estudos sobre casas compartilhadas;
  \item implicações de design para rotinas em casas compartilhadas.
\end{itemize}

O restante do artigo está organizado da seguinte forma. A
Seção~\ref{sec:fundamentacao} apresenta a fundamentação teórica e os
trabalhos relacionados. A Seção~\ref{sec:metodo} descreve o método. A
Seção~\ref{sec:resultados} apresenta os resultados da inspeção,
discutidos na Seção~\ref{sec:discussao}. A Seção~\ref{sec:limitacoes}
apresenta as limitações, e a Seção~\ref{sec:conclusao}, as conclusões e
os trabalhos futuros.

% =====================================================================
%  2 Fundamentação Teórica
% =====================================================================
\section{Fundamentação Teórica}
\label{sec:fundamentacao}

Esta seção apresenta os conceitos que sustentam a análise: a Engenharia
Semiótica e o Método de Inspeção Semiótica, os papéis que as pessoas
assumem em casas inteligentes compartilhadas, as rotinas como forma de
desenvolvimento pelo usuário final e a noção de integridade contextual.

\subsection{Engenharia Semiótica e metacomunicação}
\label{sec:engsem}

A Engenharia Semiótica é uma teoria explicativa da IHC que trata a
interface de um sistema interativo como uma comunicação enviada por seu
designer aos usuários~\cite{desouza2005semiotic}. Essa comunicação é
unidirecional: o designer não está presente durante o uso e fala por meio
da interface, que os usuários interpretam enquanto interagem com
ela~\cite{fortunato2025ihc}. Por se tratar de uma comunicação sobre como
se comunicar com o sistema, a teoria a denomina metacomunicação.

O conteúdo da metamensagem pode ser parafraseado por um modelo em que o
designer diz ao usuário qual é o seu entendimento de quem ele é, do que
ele quer ou precisa fazer, de que maneiras prefere fazê-lo e por quê; e
que, por isso, aquele é o sistema projetado para ele, que pode ou deve
ser usado de determinada forma para atingir um conjunto de objetivos
compatíveis com essa visão~\cite[p.~25]{desouza2005semiotic}. Esse modelo
se dirige a um interlocutor, o ``você'' a quem o designer fala. Em
sistemas de uso individual, esse interlocutor coincide com a pessoa que
usa o sistema. Em um dispositivo instalado em um espaço coletivo, cabe
perguntar quem é esse ``você'', questão retomada na
Seção~\ref{sec:multiusuario}.

A metamensagem é composta por signos, classificados em três
tipos~\cite{desouza2006sim,desouza2009methods}. Os \textbf{signos
estáticos} expressam o estado do sistema e podem ser interpretados sem
relações causais ou temporais, considerando apenas os elementos presentes
na interface em um dado momento, como rótulos, estados de botões e opções
selecionadas. Os \textbf{signos dinâmicos} expressam o comportamento do
sistema e só podem ser percebidos por meio da interação, como o efeito de
acionar um botão. Os \textbf{signos metalinguísticos} referem-se a outros
signos da interface e são usados pelo designer para explicar o que eles
significam e como usá-los, como mensagens de ajuda, dicas e mensagens de
erro~\cite{fortunato2025ihc}.

A qualidade de um sistema interativo, nessa perspectiva, é a sua
comunicabilidade: a capacidade de comunicar de forma organizada e
consistente a lógica, a intenção e os princípios de
design~\cite{desouza2005semiotic,desouza2009methods}. A teoria propõe dois
métodos principais para avaliá-la. O Método de Avaliação de
Comunicabilidade (MAC) observa usuários reais e avalia como a
metamensagem é recebida. O Método de Inspeção Semiótica (MIS) é conduzido
por especialistas e avalia como a metamensagem é
emitida~\cite{desouza2009methods,fortunato2025ihc}. Por não envolver
participantes, o MIS é o método adotado neste trabalho.

\subsection{O Método de Inspeção Semiótica e suas variantes}
\label{sec:mis}

No MIS, o avaliador analisa o que cada tipo de signo comunica,
reconstrói a metamensagem do designer e identifica potenciais rupturas de
comunicação~\cite{desouza2006sim}. O método é organizado nas seguintes
etapas~\cite{desouza2006sim,fortunato2025ihc}:

\begin{itemize}
  \item \textbf{Preparação:} inspeção informal, definição do foco da
        análise, dos usuários pretendidos, dos objetivos apoiados pelo
        sistema e dos cenários de inspeção;
  \item \textbf{Etapas 1, 2 e 3:} inspeção dos signos metalinguísticos,
        estáticos e dinâmicos, respectivamente, com a reconstrução de uma
        versão da metamensagem e a identificação de potenciais rupturas a
        partir de cada tipo de signo;
  \item \textbf{Etapa 4:} contraste das três versões da metamensagem, em
        busca de significados contraditórios ou ambíguos;
  \item \textbf{Etapa 5:} apreciação final da comunicabilidade do sistema,
        registrada em um relatório.
\end{itemize}

Durante a inspeção, o avaliador atua como representante dos usuários e
deve ser especialista em IHC e conhecer a Engenharia
Semiótica~\cite{desouza2006sim,fortunato2025ihc}.

Uma revisão sistemática recente identificou cinco variantes do método,
que, somadas à versão original, formam seis versões~\cite{fortunato2025ihc}.
A \textbf{MIS Científica}~\cite{desouza2010scientific} acrescenta uma
questão de pesquisa na preparação e uma etapa de triangulação dos
resultados com evidências de outras fontes, como materiais publicados na
web sobre o sistema, entre eles opiniões de usuários e perguntas
frequentes~\cite{fortunato2025ihc}. A MIS com signos dêiticos inclui uma
nova classe de signos, que se refere a elementos externos à interface: a
pessoa (identidade do usuário), o espaço e o tempo da
interação~\cite{fortunato2025ihc}. As demais variantes inspecionam mais de
uma versão da metamensagem: a CP-SIM, entre plataformas de um mesmo
sistema; a MIS mediada por leitor de tela
(SIM-SR)~\cite{carvalho2023simsr}, entre a interface original e a
interface percebida por meio do leitor de tela; e a SIM-HInt, entre os
componentes de sistemas de integração humano-computador~\cite{fortunato2025ihc}.

As variantes ampliam o método na direção de novos tipos de signos e de
múltiplas versões ou componentes de um sistema. Nenhuma das cinco, tal
como descritas por \citet{fortunato2025ihc}, tem como foco a existência de
mais de um destinatário da metamensagem, que é a situação examinada neste
trabalho.

\subsection{Casas inteligentes com múltiplos usuários}
\label{sec:multiusuario}

Diferentemente de computadores e smartphones, dispositivos de casa
inteligente instalados em espaços compartilhados tornam-se dispositivos
usados por várias pessoas e que afetam todas elas~\cite{geeng2019control}.
Os estudos sobre esses ambientes descrevem papéis distintos entre os
moradores. \citet{geeng2019control} usam o termo \textit{smart home
driver} para quem toma a iniciativa de adotar a tecnologia, sobretudo
instalando os dispositivos, e \textit{passive user} para quem participa
menos dessas decisões. \citet{koshy2021passenger} distinguem os usuários
\textit{pilotos}, que introduzem novas funcionalidades na casa, instalando
dispositivos ou criando rotinas, dos usuários \textit{passageiros}, que
usam o que já existe. Em seu levantamento, pelo menos 52\% dos
respondentes que moravam com outras pessoas viviam em casas com ao menos
um piloto e um passageiro~\cite{koshy2021passenger}.

Esses estudos convergem em alguns achados. Há uma concentração de
conhecimento, acesso e controle na pessoa que seleciona e instala os
dispositivos, que frequentemente também é a única com acesso à conta à
qual eles estão vinculados~\cite{geeng2019control}. Muitos
\textit{drivers} não consultam os demais moradores antes de instalar
dispositivos, e estes passam a depender deles quando algo deixa de
funcionar~\cite{geeng2019control}. O conhecimento sobre os dispositivos
flui dos pilotos para os passageiros, e a casa tende a refletir os valores
dos pilotos, o que dificulta que os passageiros incorporem os dispositivos
ao seu cotidiano~\cite{koshy2021passenger}. Em um dos domicílios
estudados por \citet{koshy2021passenger}, o piloto praticamente deixou de
criar rotinas porque elas incomodavam a esposa: as rotinas matinais, com
horário fixo e pensadas para a agenda dele, não combinavam com o horário
de trabalho dela. \citet{zeng2017concerns} observaram
um descompasso semelhante: os moradores que não administram a casa
tendem a estar menos informados sobre os riscos de segurança e
privacidade. \citet{zeng2019multiuser} mostram que esse desequilíbrio de
poder aparece também em relações não abusivas e propõem como princípios
de design a flexibilidade no controle de acesso, a agência de todos os
usuários, o respeito entre eles e a transparência do comportamento da
casa.

Além dos moradores, há pessoas que convivem com os dispositivos sem
participar de sua configuração. \citet{bernd2020bystanders} estudaram
babás que trabalham em casas com dispositivos de vigilância e as
analisaram como pessoas afetadas (\textit{bystanders}), em uma relação
marcada pela diferença de poder entre empregador e empregada.
\citet{geeng2019control} relatam uma tensão semelhante envolvendo o acesso
de uma diarista a uma fechadura inteligente.

Os assistentes de voz acentuam o caráter coletivo desses dispositivos. O
uso da Amazon Echo se entrelaça com as conversas e as atividades
simultâneas dos moradores, como as refeições em
família~\cite{porcheron2018voice}, e a interface de voz, acessível a todos,
amplia o acesso da família à tecnologia~\cite{beneteau2020parenting}. A
configuração do dispositivo e de suas rotinas, no entanto, continua
concentrada no aplicativo vinculado à conta de uma
pessoa~\cite{barricelli2024routines}.

\subsection{Rotinas como desenvolvimento pelo usuário final}
\label{sec:rotinas}

Rotinas são sequências de ações definidas pelo usuário e disparadas por
eventos ou condições específicas~\cite{barricelli2022multimodal}. Na
Amazon Alexa e no Google Nest, elas são criadas por meio do aplicativo
complementar no smartphone~\cite{barricelli2024routines}. Essa forma de
configuração corresponde à programação gatilho-ação (``se \ldots{},
então \ldots{}''), na qual \citet{ur2014tap} mostraram que é possível
expressar a maior parte dos comportamentos desejados pelos participantes
de seu estudo e que usuários sem experiência aprendem rapidamente a criar
regras com múltiplos gatilhos ou ações.

Criar a rotina, porém, é só uma parte do uso. \citet{koshy2021passenger},
retomando o estudo longitudinal de \citet{woo2015diy}, relatam que os
usuários dedicam tempo a refinar rotinas depois de criadas, abandonam
muitas delas e que os autores sugeriram oferecer um canal para que os
moradores avaliem as rotinas configuradas por outros membros da casa.
Antes da popularização dos assistentes de voz, \citet{brush2011wild} já
apontavam a pouca flexibilidade e a dificuldade de gerenciar a automação
residencial como barreiras à sua adoção.

\subsection{Integridade contextual}
\label{sec:integridade}

Para analisar o uso de dados nas rotinas, adotamos a noção de integridade
contextual. \citet{nissenbaum2004privacy} propõe que a proteção da
privacidade depende das normas de cada contexto: a coleta e a circulação
de informações devem ser adequadas ao contexto e obedecer às normas de
distribuição que o regem. A autora distingue as normas de adequação, que
definem quais informações sobre uma pessoa é apropriado revelar em um
contexto, das normas de fluxo ou distribuição, que definem como essas
informações podem circular, e afirma que não há esfera da vida em que
``tudo é permitido''. \citet{bernd2020bystanders} usam esse referencial
para analisar como as relações de poder entre moradores e pessoas afetadas
moldam as normas de privacidade em casas inteligentes. No caso das
rotinas, um dado associado a uma pessoa, como o horário de seu alarme,
pode disparar ações percebidas por todos que estão no ambiente, o que
caracteriza uma mudança no fluxo dessa informação.

% =====================================================================
%  Método
% =====================================================================
\section{Método}
\label{sec:metodo}

\pendente{escopo e versão do aplicativo; construção das duas personas a partir da literatura (conceito de persona com Pruitt e Grudin 2003 e risco de estereótipos com Marsden e Haag 2016); cenários de inspeção; etapas da MIS aplicadas a cada persona; contraste entre o que a interface comunica a cada uma}

% =====================================================================
%  Resultados
% =====================================================================
\section{Resultados}
\label{sec:resultados}

\pendente{metamensagens por classe de signo; contraste e apreciação final; rupturas encontradas para cada persona}

% =====================================================================
%  Discussão
% =====================================================================
\section{Discussão}
\label{sec:discussao}

\pendente{rupturas em relação aos GCs (GC4, GC1, GC7); implicações de design para casas compartilhadas}

% =====================================================================
%  Limitações
% =====================================================================
\section{Limitações}
\label{sec:limitacoes}

\pendente{inspeção não é uso real; personas são representações fictícias; versão e idioma do aplicativo}

% =====================================================================
%  Conclusão e Trabalhos Futuros
% =====================================================================
\section{Conclusão e Trabalhos Futuros}
\label{sec:conclusao}

\pendente{síntese das respostas às QPs; interface adaptativa para autoria coletiva de rotinas como trabalho futuro}

% =====================================================================
%  Considerações Éticas (seção obrigatória no IHC 2026)
% =====================================================================
\section{Considerações Éticas}
\label{sec:etica}

\pendente{redigir após definir o método. Base: o estudo não envolveu
participantes; as personas são fictícias e a inspeção usou contas dos
próprios pesquisadores, em conformidade com a Resolução CNS n.º 510/2016.}

% =====================================================================
%  Agradecimentos. O IHC 2026 pede a declaração explícita do uso de
%  IA generativa na escrita ou revisão, de preferência nesta seção.
% =====================================================================
\section*{Agradecimentos}

\pendente{agradecimentos e declaração do uso de ferramentas de IA
generativa na escrita e revisão do texto}


% ---------------------------------------------------------------------
%  Referências (estilo oficial SBC: citações autor-ano)
% ---------------------------------------------------------------------
\bibliographystyle{sbc}
\bibliography{references}

\end{document}
